\section{Phase 1 and the stock bottleneck map}
\label{sec:baseline}

\subsection{Reproducible baseline}
Phase~1 established a reproducible reference: \code{MACECalculator(device="cuda")} on ASE's H$_2$O, one warm-up and ten timed \code{reset}$\to$\code{get\_potential\_energy}$\to$\code{get\_forces} calls, CUDA-synchronised. The calculator leaves torch's default dtype at FP32 and silently switches to FP64 to match the checkpoint; e3nn 0.4.4's \code{torch.load} of its constants needed \code{TORCH\_FORCE\_NO\_WEIGHTS\_ONLY\_LOAD=1} under torch $\ge2.6$. Energies ($-2081.11639716893$\,eV) and forces were bitwise identical across runs and across a full rebuild of the environment. The H$_2$O call took 14.2--14.7\,ms (median), of which almost nothing is GPU work. One bookkeeping caveat: the Phase-1 JSON files were later overwritten by a post-build rerun; the original numbers survive in \code{logs/04} (14.31/14.29\,ms) and agree within 2\%.

\subsection{Where stock MACE spends its time}
Phases 3 and 6 decomposed the stock call. \Cref{tab:stockprof} gives the Phase-6 profile (job 1017, torch.profiler/CUPTI, graph built once and kept on GPU so that only the model is measured).

\begin{table}[h]
\centering
\caption{Eager model phase totals (ms) on a prebuilt graph, job 1017. ``busy'' is the union of kernel intervals; idle $=1-\text{busy}/\text{wall}$; nK = kernels launched. flash = FlashTP conv + e3nn SC; hybrid = FlashTP conv + cuEquivariance SC.}
\label{tab:stockprof}
\small
\begin{tabular}{llrrrrrr}
\toprule
Backend & System & fwd wall & fwd busy & idle & nK fwd & fwd+bwd wall & nK f+b \\
\midrule
e3nn (stock) & aspirin & 6.11 & 0.87 & 86\% & 295 & 13.76 & 830 \\
flash & aspirin & 5.81 & 0.75 & 87\% & 252 & 11.46 & 616 \\
hybrid & aspirin & 8.66 & 0.49 & 94\% & 221 & 17.01 & 556 \\
e3nn (stock) & water648 & 14.17 & 12.60 & 11\% & 304 & 36.63 & 864 \\
flash & water648 & 9.93 & 8.12 & 18\% & 256 & 21.12 & 637 \\
hybrid & water648 & 8.71 & 2.49 & 71\% & 219 & 17.35 & 559 \\
\bottomrule
\end{tabular}
\end{table}

Three distinct bottlenecks emerge:

\paragraph{(1) A host launch floor.} A forward issues 250--300 kernels and a force backward 330--560. At aspirin the GPU is busy for 0.9\,ms of a 6.1\,ms forward. Only 10--15\% of the CPU time is in \code{cudaLaunchKernel}; the rest is Python, the dispatcher, and library host code. The cuEquivariance symmetric contraction is the extreme case: 4.6\,ms of profiled host time per forward (mostly \code{cuequivariance::uniform\_1d} host calls) for 15 kernels of 37--70\us, and another 3.6\,ms in backward. A segment micro-benchmark (Phase 6, job 1019) found the cueq product block costs 8.3\,ms of host time for 0.12--0.24\,ms of GPU work. This floor makes the hybrid backend \emph{slower} than stock below $\sim$400 atoms when run eagerly.

\paragraph{(2) Generic equivariant operators.} At water648 the e3nn convolution costs 5.7\,ms forward and 13.6\,ms backward, dominated by ATen elementwise kernels on $[E,2048/5120]$ FP64 tensors (multiply 17\%, add 14\%, copy 9\% of GPU time) plus FP64 GEMMs (14\%); the e3nn symmetric contraction adds 5.7/6.6\,ms. Together they are $\sim$85\% of GPU time. Memory saved for backward is 4.24\,GB for e3nn (2.9\,GB of it conv edge tensors) versus 0.76\,GB for the hybrid.

\paragraph{(3) Synchronous host-side graph construction.} \Cref{tab:graphbuild} decomposes MACE's \code{\_atoms\_to\_batch} (Phase 5). matscipy's neighbour list is only part of the cost; the rest is $\sim$30 small \code{torch.tensor} constructions, a collate, and 26--37 pageable host-to-device copies, each followed by a stream synchronisation.

\begin{table}[h]
\centering
\caption{Stock graph construction, median ms (Phase 5, job 1005).}
\label{tab:graphbuild}
\small
\begin{tabular}{lrrrrr}
\toprule
Part & H$_2$O & water192 & water648 & water1536 & wbox648 \\
\midrule
matscipy neighbour list & 0.037 & 0.717 & 3.472 & 9.374 & 7.049 \\
MACE NL wrapper (shifts, self-edges) & 0.103 & 0.317 & 1.037 & 2.290 & 1.248 \\
from\_config other & 0.215 & 0.277 & 0.341 & 0.432 & 0.374 \\
collate & 0.741 & 0.855 & 1.108 & 1.205 & 1.135 \\
copy to GPU & 0.910 & 1.049 & 1.345 & 1.576 & 1.397 \\
clone on GPU & 0.328 & 0.356 & 0.454 & 0.482 & 0.468 \\
\midrule
\textbf{total} & 2.581 & 3.892 & 8.283 & 16.087 & 12.182 \\
\bottomrule
\end{tabular}
\end{table}

\subsection{Consequence for the plan}
The three bottlenecks dominate in different regimes, so no single fix could help across sizes: a faster convolution does nothing for aspirin, and removing the launch floor does little for 5{,}000 atoms. Every subsequent phase attacked one of them, in roughly the order their share grew as the previous one shrank. \Cref{fig:ladder} (\cref{sec:results}) shows the resulting ladder.
